% ======================================================================
% PLOS ONE submission -- prepared from the official PLOS LaTeX template
% (Version 3.8, April 2026). Bibliography style: plos2025.bst
% ======================================================================
\documentclass[10pt,letterpaper]{article}
\usepackage[top=0.85in,left=2.75in,footskip=0.75in]{geometry}
\usepackage{tabularx}
\usepackage{seqsplit}
\usepackage{booktabs}
\usepackage{amsmath,amssymb}
\usepackage{changepage}
\usepackage{textcomp,marvosym}
\usepackage{cite}
\usepackage{nameref,hyperref}
\usepackage[right]{lineno}
\usepackage[nopatch=eqnum]{microtype}
\microtypesetup{expansion=false}
\DisableLigatures[f]{encoding = *, family = * }
\usepackage[table]{xcolor}
\usepackage{array}

% --- packages required by this manuscript's tables/figures ---
\usepackage[T1]{fontenc}
\usepackage{float}
\usepackage{multirow}
\usepackage{booktabs}

% create "+" rule type for thick vertical lines
\newcolumntype{+}{!{\vrule width 2pt}}
\newlength\savedwidth
\newcommand\thickcline[1]{%
  \noalign{\global\savedwidth\arrayrulewidth\global\arrayrulewidth 2pt}%
  \cline{#1}%
  \noalign{\vskip\arrayrulewidth}%
  \noalign{\global\arrayrulewidth\savedwidth}%
}
\newcommand\thickhline{\noalign{\global\savedwidth\arrayrulewidth\global\arrayrulewidth 2pt}%
\hline
\noalign{\global\arrayrulewidth\savedwidth}}

% Remove comment for double spacing
\usepackage{setspace}
\doublespacing

% Text layout
\raggedright
\setlength{\parindent}{0.5cm}
\textwidth 5.25in
\textheight 8.75in

% Bold 'Fig #' in caption, period separator, left justified
\usepackage[aboveskip=1pt,labelfont=bf,labelsep=period,justification=raggedright,singlelinecheck=off]{caption}
\renewcommand{\figurename}{Fig}

% Please use the included `plos2025.bst` as the BibTeX style
\bibliographystyle{../../plos2025}

% Remove brackets from numbering in List of References
\makeatletter
\renewcommand{\@biblabel}[1]{\quad#1.}
\makeatother

% Header and Footer with logo
\usepackage{lastpage,fancyhdr,graphicx}
\usepackage{epstopdf}
\pagestyle{fancy}
\fancyhf{}
\rfoot{\thepage/\pageref{LastPage}}
\renewcommand{\headrulewidth}{0pt}
\renewcommand{\footrule}{\hrule height 2pt \vspace{2mm}}
\fancyheadoffset[L]{2.25in}
\fancyfootoffset[L]{2.25in}
\lfoot{\today}

\begin{document}
\vspace*{0.2in}

% Title must be 250 characters or less.
\begin{flushleft}
{\Large
\textbf{Coevolutionary balance of resting-state brain networks in autism}
}
\newline
\\
S. Rezaei Afshar\textsuperscript{1},
G. R. Jafari\textsuperscript{2,*}
\\
\bigskip
\textbf{1} Institute for Cognitive and Brain Sciences, Shahid Beheshti University, Evin, Tehran, Iran
\\
\textbf{2} Department of Physics, Shahid Beheshti University, Evin, Tehran, Iran
\\
\bigskip

* Corresponding author: gjafari@gmail.com

\end{flushleft}

\noindent\textbf{Working revision draft: author review required.}
This document replaces the submitted analysis with the verified revision outputs.
Historical fALFF provenance and the public archive identifier remain unresolved
and are explicitly marked below.

\section*{Abstract}
Signed functional connectivity and regional low-frequency BOLD fluctuations offer
complementary descriptions of resting-state organization. We examined their joint
alignment using a static Hamiltonian inspired by coevolutionary balance in a
male-only adult sample from ABIDE I. Downloaded preprocessed regional time series
and fALFF maps were analyzed with and without global signal regression (GSR).
Age and IQ eligibility yielded 140 participants (69 autistic participants and
71 controls). Measurement coverage defined a fixed common set of 106 regions;
site-adjusted group models included 136 participants. In a separate male-control
baseline, mean empirical Hamiltonian values were lower than their null expectations
under unrestricted and network/hemisphere-blocked state permutations. Hamiltonian values did not
show a detectable diagnosis effect after adjustment for site, age, full-scale IQ,
and mean framewise displacement (GSR: coefficient 1.43, 95\% CI $-17.17$ to 20.04;
no-GSR: 2.06, 95\% CI $-12.39$ to 16.50; both Holm-adjusted $p=1.00$).
ComBat, pooled energy-trimming, and larger-node-set sensitivities did not identify
a corrected diagnosis effect. Equivalence tests were inconclusive. Exploratory
site-held-out classification yielded balanced accuracies of 0.457--0.558, with
all conditional confidence intervals including 0.5.
A broader exploratory selection-method search did not establish a reliable performance gain. Clinical analyses and
motion- and mean-connectivity-adjusted efficiency analyses provided no corrected
evidence of an association. These results do not establish group equivalence,
diagnostic utility, or an autism-specific energy--efficiency relationship.
Interpretation is restricted to the selected adult males, the coverage-limited
network, and the downloaded derivatives; historical fALFF processing provenance
must be resolved before final submission.

\clearpage
\newgeometry{top=0.85in,left=1in,right=1in,footskip=0.75in}
\linenumbers

\section*{Introduction}
Resting-state functional MRI provides several descriptions of brain organization.
Functional connectivity (FC) measures statistical dependence between regional
BOLD time series, whereas fractional amplitude of low-frequency fluctuations
(fALFF) describes the relative contribution of low-frequency fluctuations to the
sampled spectrum \cite{Zou:2008fa}. Neither quantity directly measures neuronal
firing rate. Regional fALFF differences have been reported in adult autism
\cite{Itahashi2015revision}, but such regional effects need not produce a difference
in a scalar summary combining fALFF ranks and FC.

Signed networks retain positive and negative associations, whose treatment can
materially affect network summaries \cite{RubinovSporns2011}. Structural balance
theory describes consistency among signed relationships \cite{Heider:1946ac,cartwright1956}.
Coevolutionary models introduce nodal states as well as signed interactions
\cite{Kargaran:2021he,GhanbarzadehNoudehi2022}. Here, these ideas motivate a static
alignment statistic. We do not estimate a dynamical coevolutionary process,
metabolic energy, or excitation--inhibition balance from BOLD data.

Our earlier ABIDE structural-balance analysis studied triadic relationships in
children and adolescents aged 6--18 without GSR \cite{Moradimanesh2021}. Its
statistic and developmental population differ from the present adult nodal-state
Hamiltonian. Likewise, energy-landscape analyses of autism model time-varying
state occupancy and transitions \cite{Watanabe2017revision}; a static alignment
statistic cannot reproduce or refute a finding about those transitions.

Multisite FC studies show substantial heterogeneity and limited generalization
across acquisition sites \cite{King2019revision}. Site adjustment and harmonization
therefore require explicit treatment, particularly when diagnosis and site are
unbalanced. Protecting diagnosis in a harmonization model does not guarantee
unbiased downstream inference \cite{Nygaard2016revision}. We use raw adjusted
features for primary inference and treat ComBat as a sensitivity analysis.

This revision asks whether a coverage-defined alignment statistic differs by
diagnosis in adult males with observed full-scale IQ (FIQ) above 80. We compare
GSR and no-GSR derivatives, quantify uncertainty, examine the role of mean FC and
motion in efficiency associations, and assess exploratory site-held-out prediction.
The revised node and participant definitions were fixed using demographic criteria
and measurement QC before corrected group-effect testing; they are revision-stage
choices, not an original preregistration.

\section*{Results}
\subsection*{Participant and measurement scope}
Figure~\ref{fig:flow} separates demographic eligibility, feature availability, and
site-adjusted inference. Table~\ref{tab:demographics} reports the 136 modeled
participants. The nine-site model includes an ASD-only site, CALTECH ($n=9$);
its diagnosis effect is identified through the remaining mixed-diagnosis sites,
under the common-effect regression assumption. Broader-181 and full-200 designs
have 101 (47/54) and 64 (30/34) adjusted-model participants, respectively.

\input{tables/demographics.tex}

\subsection*{Null-model comparisons of coevolutionary alignment}
\paragraph*{Male-control baseline.}
We first examined the 71 eligible male controls as a separate non-autistic
baseline, rather than pooling them with ASD. Each control contributes 1,000
state permutations per pipeline and null scheme. Mean empirical $H$ is $-91.39$
with GSR and $-66.62$ without GSR. The corresponding unrestricted null means
are 1.49 and 17.11; mean empirical-minus-null differences are $-92.88$ and
$-83.73$, with descriptive site-bootstrap 95\% intervals of $-103.90$ to
$-79.89$ and $-89.79$ to $-77.66$, respectively.

With network/hemisphere-constrained states, the null means are $-41.35$ and
$-18.85$, and the empirical-minus-null differences are $-50.04$ (95\% interval
$-54.14$ to $-44.03$) and $-47.76$ ($-53.11$ to $-42.16$).
Table~\ref{tab:controlnull} and Fig.~\ref{fig:controlnull} show these comparisons;
individual null $z$-scores are plotted in S1 Supporting Information.
All 71 controls in each pipeline have empirical $H$ below their unrestricted
null mean. Under the blocked null, this holds for 69/71 with GSR and 67/71
without GSR. After exploratory BH correction over all 284 individual control
comparisons, the numbers with lower-tail $q<.05$ are 61/71 and 64/71 for the
unrestricted null and 36/71 and 50/71 for the blocked null (GSR/no-GSR).

Thus empirical control state assignments tend to align with their own FC more
strongly than the randomized assignments under both specified nulls, with
substantial dependence on the constraints. This is evidence about the chosen
alignment statistic and randomizations, not proof of dynamical equilibrium,
a comparison with topology-rewired graphs, or an autism-specific effect.

\input{tables/control_null_summary.tex}

\paragraph*{ASD and across-diagnosis summaries.}
All 140 participants per pipeline have finite null scores under both randomizations.
Within ASD alone ($n=69$), mean unrestricted $z$ is $-4.705$/$-4.917$
(GSR/no-GSR), and blocked mean $z$ is $-2.331$/$-2.591$. With exploratory
BH across all 560 participant/pipeline/null comparisons, the ASD lower-tail
$q<.05$ counts are 61/69 and 60/69 for unrestricted states, and 40/69 and
46/69 for blocked states. These complements retain the original ASD as well
as control null question and do not test diagnosis differences.
The corresponding across-diagnosis summaries are retained for completeness.
Mean unrestricted $z$ is $-4.70$/$-5.04$ (GSR/no-GSR), compared with
$-2.32$/$-2.68$ under network/hemisphere blocks. These values show dependence on
the randomization scheme, not diagnosis discrimination or preservation of full
spatial structure. Individual scores and probabilities are in S2 Supporting Information.



\subsection*{Whole-network coevolutionary energy and motif composition}
Primary adjusted coefficients are 1.432 with GSR (95\% CI $-17.173$ to 20.038,
$p=.879$) and 2.056 without GSR (95\% CI $-12.387$ to 16.499, $p=.779$);
both Holm-adjusted $p=1.000$ (Table~\ref{tab:primary}, Fig.~\ref{fig:primary}).
No valid secondary group test survives its specified correction in any raw or
ComBat design. The trimming sensitivities likewise do not identify a corrected
$H$ difference. Larger-node results are in S1 Supporting Information.

\input{tables/primary_H.tex}

Figure~\ref{fig:globalfeatures} displays all seven valid global quantities for
the untrimmed 140-person feature cohort. Corrected within- and between-network
effect heatmaps in S1 show all prespecified cortical blocks rather than selecting
nominally significant examples.

All six ComBat design/pipeline pairs complete with finite, changed feature values.
Primary descriptive site partial $R^2$ decreases from .056 to .004 with GSR and
from .047 to .005 without GSR. Maximum primary leverage is .352; maximum Cook's
distance is .080 or smaller across these four raw/ComBat fits. These checks do
not establish perfect harmonization or independence of participants within sites.
Excluding CALTECH yields 127 participants (58 ASD, 69 controls) at eight sites;
raw adjusted coefficients remain imprecise (GSR: 1.115, 95\% CI $-17.531$ to
19.762; no-GSR: 1.660, 95\% CI $-12.775$ to 16.096; both Holm $p=1.000$).

Primary unadjusted Hedges' $g$ is $-.024$ (bootstrap 95\% CI $-.384$ to .315)
with GSR and .072 ($-.268$ to .412) without GSR. Approximate two-sample 80\%-power
minimum detectable $d$ is .484 at $\alpha=.05$, or .534 at $\alpha=.025$.
Equivalence tests are inconclusive ($p=.193$ and .223). Nonsignificant contrasts
therefore do not demonstrate absence of a meaningful diagnosis effect.

\subsection*{Network-specific coevolutionary metrics}
We retain the original within-network question: whether ASD and control
participants differ in mean FC (AvgWeight), mean state-weighted FC (AvgFalffW),
or block Hamiltonian contribution (EnergyBlock) within each of the seven Yeo
systems. The primary graph permits 18 tests per pipeline; three within-limbic
quantities are unavailable because only one limbic node remains. All available
blocks are shown in Fig.~\ref{fig:intraeffects}, with complete estimates in S2.
No comparison survives BH across the 36 pipeline/within-network tests. In GSR,
the smallest raw probability is .073. Without GSR, somatomotor mean FC is
nominally lower in ASD (coefficient $-.0683$, $p=.0371$), as is default-network
mean FC ($-.0519$, $p=.0432$); both become $q=.392$ and are not corrected findings.

\subsection*{Inter-network connectivity and energy}
We also retain all 21 unordered network pairs and all three original measures,
yielding 63 tests per pipeline. Fig.~\ref{fig:intereffects} shows every block,
including nonsignificant ones. No comparison survives BH across the 126
pipeline/inter-network tests. GSR salience/ventral-attention--limbic state-weighted
FC and its corresponding block energy are nominally different ($p=.0315$,
$q=.756$). They are proportional descriptions of the same block, not two
independent effects. No no-GSR inter-network test reaches $p<.05$; its smallest
raw probability is .0508. These results address spatially specific diagnosis
effects separately from the global Hamiltonian and the within-group associations.

\subsection*{Relationships between H and other network metrics}
To preserve the original descriptive analysis, we correlate $H$ with all 87 other
available nonlegacy features within ASD and controls separately, in both pipelines.
All eligible participants are included (69 ASD/71 controls). Under the original
per-group Pearson BH convention, 26 GSR-ASD, 21 GSR-control, and 12 no-GSR-ASD
associations survive $q<.05$; none survives in no-GSR controls
(Table~\ref{tab:Hcorrelations}). Figure~\ref{fig:Hcorrelates} restores the
original top-20 presentation, while S2 contains every test, not just ranked results.

Examples include GSR-ASD $H$ versus somatomotor--default block energy
($r=.434$, $q=.00472$) and no-GSR-ASD $H$ versus frontoparietal--default block
energy ($r=.473$, $q=.000890$). EnergyBlock and AvgFalffW are proportional within
a fixed block, and both share the state-weighted terms entering $H$; these
associations are therefore not independent validation. Among mean-FC features,
GSR controls show an association with somatomotor--default AvgWeight
($r=-.315$, $q=.0346$), still derived from the same FC matrix. Significance in
one group alone does not establish a group difference in correlation.

The descriptive raw $H$--global-efficiency correlations are .022 in GSR-ASD and
.109 in no-GSR-ASD, neither significant; corresponding control correlations are
.100 and .140. The adjusted analysis below addresses the additional question of
association conditional on mean FC, motion, age, IQ, and site. Matched-participant
raw/ComBat comparisons and alternative pooled corrections are retained in S2,
without selecting the version with the most significant associations.

\input{tables/H_correlation_counts.tex}

\subsection*{H--efficiency relationships after nuisance adjustment}

ASD partial $H$--global-efficiency correlations are .017 with GSR and .077 without
GSR ($n=66$); neither is significant after correction. Diagnosis-by-efficiency
interactions are also nonsignificant (Holm $p=.635$ and .940), providing no
support for an ASD-specific trade-off. Residualized efficiency scatter plots and
complete clinical validation plots are provided in S1 Supporting Information.

\subsection*{Correlation analysis of network metrics with behavioral scores}
All five original ADI-R/ADOS outcomes are retained with score-specific missingness
and a common 88-feature schema across pipelines. Pearson and Spearman analyses
are both reported. No clinical feature correlation survives
FDR; adjusted ADOS-total correlations with $H$ are .008 and .061 ($n=50$).

\paragraph*{Multivariate analysis: Partial Least Squares.}
PLSR and PLSC remain distinct complements, as in the original paper.
Score-specific PLSR samples range from 28 to 50; every OOF $R^2$ is negative and
every MSE exceeds its training-mean baseline. No PLSC test survives BH correction;
a nominal no-GSR five-score LV3 probability of .007 becomes $q=.098$ and is not
presented as a confirmed latent association.

\subsection*{Machine learning classification of ASD versus controls}
% CLASSIFICATION_RESULTS_BEGIN
The 10 completed classifier/pipeline pairs have 140 unique OOF predictions
each, using the same frozen site-disjoint splits. Completed algorithms are
logistic regression, RBF SVM, Gaussian naive Bayes, KNN, XGBoost. Balanced accuracy ranges from
0.457 to 0.558 (Table~\ref{tab:classification}, Fig.~\ref{fig:primary}).
Every conditional interval includes .5. No final classifier is selected from outer-fold performance.
The conditional intervals do not include all training or QC-selection uncertainty.
% CLASSIFICATION_RESULTS_END

\input{tables/classification.tex}

% ML_SELECTION_RESULTS_BEGIN
\paragraph*{Feature-selection audit and exploratory extension.}
Feature selection was already nested in the five-model analysis above. Reconstructing
all 50 outer-training fits reproduces every saved prediction: ANOVA retained 22
features in 11 fits, 44 in 17, and all 88 in 22. We subsequently compared broader
selection methods in a separate exploratory analysis, retaining the same known
outer site folds rather than treating them as a fresh independent test.
Training-fold near-exact redundancy pruning left 59 distinct columns, and inner
CV chose all 59 in seven fits, ANOVA top-K in 16, mutual-information top-K in 17,
and L1-ranked top-K in ten. Additional selection retained 5, 10 or 20 features.

Four of the ten model/pipeline point estimates increased and six decreased
(Table~\ref{tab:FScomparison}). No-GSR logistic regression changed from .485 to
.544, no-GSR SVM from .485 to .536, and no-GSR KNN from .542 to .557. GSR XGBoost
changed from .457 to .492. Every paired conditional interval for change includes
zero, and all follow-up balanced-accuracy intervals include .5. The follow-up
therefore does not establish a reliable performance gain or justify selecting a
new final classifier from the outer results. Full ROC, confusion, performance
and selected-feature frequency heatmaps are in S1 Supporting Information; all
feature names, inner candidates, and test predictions are in S2.

\input{tables/feature_selection_comparison.tex}
% ML_SELECTION_RESULTS_END

\subsection*{Sensitivity analysis: results without global signal regression}
The original no-GSR sensitivity remains explicit at each scale: global quantities,
within-network and between-network group contrasts, H--metric associations,
clinical analyses, PLS, and classification. Primary participant and node sets are
identical across pipelines. GSR/no-GSR Hamiltonian diagnosis estimates are both
imprecise; ComBat and pooled trimming do not yield a corrected group difference.
Larger-node sensitivity designs change participants as well as nodes.
The original spectral-bipolarity descriptor is now reported separately as an
exploratory raw adjusted comparison: coefficients are .00040 with GSR and .01055
without GSR, with Holm $p=.916$ for both. It remains outside the frozen feature
families and does not alter the common clinical/predictive schema.

% CLASSIFICATION_CONTINUITY_BEGIN
The revised site-held-out analysis now restores 5 of the five original
algorithm types, using the same revision participants, features, folds and scoring.
All five original algorithm types have complete OOF predictions in both pipelines.
Matching the algorithm types does not restore identical original inputs or validation:
site-held-out evaluation replaces the original random train/test procedure.
% CLASSIFICATION_CONTINUITY_END


\section*{Discussion}
The revised analysis does not detect a corrected diagnosis effect in the
coverage-defined adult male sample. Raw covariate-adjusted models, explicit
ComBat sensitivities, pooled trimming, and mixed-diagnosis-site restriction are
consistent in this limited sense. The separate male-control null comparison supports alignment relative to the
specified state randomizations, while the diagnosis comparison asks a different
question. Alignment in controls can coexist with no detectable diagnosis contrast.
Their uncertainty remains substantial, and
inconclusive equivalence tests preclude a universal adult-null interpretation.
The multisite generalization difficulties reported for FC more broadly
\cite{King2019revision} reinforce the need for larger independent tests.

This result concerns a relative fALFF--FC alignment statistic. Median states do
not distinguish globally elevated from globally reduced fALFF distributions,
and identical values can receive different labels across participants. It is
therefore inappropriate to equate $+1$ with high neuronal firing or to interpret
$H$ as metabolic cost. Construction also links the statistic to global weights;
the original descriptive H--metric question still produces corrected
within-group associations, whereas the controlled efficiency analyses do not
retain the submitted ASD-specific trade-off interpretation. An association between two summaries of the same FC matrix would
not independently validate a biological mechanism even if statistically detected.

The result differs in scope from our earlier developmental triadic-balance work
\cite{Moradimanesh2021} and from temporal energy-landscape studies
\cite{Watanabe2017revision}. Different ages, state definitions, parcellations, and
statistics can test different hypotheses. Regional adult fALFF findings
\cite{Itahashi2015revision} can coexist with a null global rank-based summary.
We cannot infer developmental normalization or disagreement with dynamic-state
findings from this adult cross-sectional analysis.

\subsection*{Sex and generalizability}
The male-only design restricts both the population represented and the questions
that can be answered. The selected ABIDE pool does not support an inference about
autistic women, girls, nonbinary people, or people with intellectual disability.
The practical scarcity of eligible female participants in a dataset does not
justify extending male estimates to all autistic people. Repeated male-only
sampling can leave diagnostic frameworks and candidate biomarkers poorly tested
in underrepresented groups and may contribute to their under-recognition.

Sex-dependent diagnostic FC patterns have been reported \cite{Alaerts2016revision}.
Sex-related variation in either FC or fALFF could affect their joint alignment,
even when within-person median splitting removes absolute fALFF scaling.
Consequently, neither this group-null result nor the absence of a corrected
efficiency association can be assumed to generalize to autistic females.
Recorded binary sex metadata also do not capture gender identity or the social
and clinical pathways influencing ascertainment. An inclusive test should recruit
sufficient participants across sexes, ages, cognitive abilities, and presentation
profiles; prespecify sex-by-diagnosis analyses and measurement-QC comparisons;
and validate calibration and performance in each group. A small female subgroup
added only for an underpowered post hoc comparison would not resolve this gap.

\subsection*{Other limitations and revision transparency}
Coverage selection removes 94 CC200 nodes from the primary graph and leaves only
one limbic node; the result is not a complete whole-brain or limbic test. The
50\% support rule can still permit heterogeneous spatial sampling within a parcel.
Larger-node designs trade coverage for participant retention and do not isolate
a node-count effect. Signed associations depend on nuisance processing, and
unmodeled site/acquisition effects may remain. HC3 addresses leverage and
heteroscedasticity at participant level; it is not site-clustered uncertainty.
Clinical samples are small, permutation assumptions are approximate, and coarse
block nulls do not preserve full spatial autocorrelation
\cite{Markello2021revision}. Prediction is conditional on cohort-wide QC.

The original within-ASD no-GSR H--efficiency association can be reproduced from
the supplied original feature exports ($r=-.324$, $q=.00592$, $n=80$). Its absence
in the revision's 106-node definition does not by itself establish that the
original association was false: the participants, nodes, fALFF inputs and model
assumptions differ. Nor does reproduction of that association establish a
between-diagnosis interaction or independent biological mechanism. We preserve
both the original record and the revised analyses, stating their different scopes.

The original-code audit identified discrepancies in cohort flow, ordinal cortical
assignments, a fixed-TR spectral ratio calculated from filtered ROI series, and
exports labeled ComBat that equaled raw values after a failed harmonization path.
The revised analysis uses downloaded fALFF maps, explicit anatomical assignments,
and checked harmonization. Multiple elements changed together, so we cannot
attribute changed results solely to correction of site effects.
Historical PCP fALFF input/band provenance remains a submission blocker; this
working draft does not establish that preprocessing concern has been resolved.

\paragraph*{Conclusions.}
Within the specified coverage-limited network of adult males, the revised analyses
provide no corrected evidence of diagnosis differences or clinical/predictive
utility. Male controls show alignment relative to the specified state-permutation
nulls, a baseline property distinct from diagnosis discrimination. These analyses
do not establish group equivalence, a universal autism null, or
an ASD-specific energy--efficiency mechanism. Independent, inclusive validation
and verified derivative provenance are necessary for stronger conclusions.

\section*{Methods}
\subsection*{Data source, eligibility, and ethics}
The source pool comprises 182 previously selected ABIDE I participants (90 ASD,
92 controls) with requested PCP derivatives \cite{DiMartino:2014abide,PCPdocumentationRevision}.
We did not reconstruct an initial random or matched 200-person sample. Eligibility
for the revision required age 18--30 inclusive and observed FIQ $>80$, without IQ
imputation, yielding 140 participants (69 ASD, 71 controls) at 12 sites. All were
recorded male. We did not impose a right-handedness restriction: 71 were recorded
right-handed, six left-handed, three ambidextrous, and 60 missing or sentinel-coded.
Existing diagnoses were taken from ABIDE metadata; we do not retrospectively assign
DSM-5 ascertainment to all contributing sites. Clinical missingness was handled
per outcome rather than excluding everyone missing any clinical score.

The sample flow is shown in Fig.~\ref{fig:flow}. Model-specific counts and
covariates are given in Table~\ref{tab:demographics}; site counts and participant
ledgers are supplied separately in S1 and S2 Supporting Information.

The original ABIDE studies were approved by the contributing sites' institutional
review boards, with written informed consent and/or assent obtained under their
respective protocols \cite{DiMartino:2014abide}. The present study exclusively
reanalyzed publicly available, deidentified ABIDE data. No new participants were
recruited, no identifying information was accessed, and there was no participant
contact or new data collection. No additional ethics approval or consent was
obtained for this secondary analysis because it used only those public,
deidentified records.

\subsection*{Downloaded derivatives and upstream processing}
Regional CC200 time series, fALFF maps, and functional masks were downloaded from
the PCP C-PAC outputs. MRI preprocessing, time-series extraction, and voxelwise
fALFF computation were performed upstream, not locally for this revision.
The PCP documentation describes slice-time correction, motion correction,
intensity normalization, and nuisance regression using six motion parameters,
their one-time-point lags, the squares of those 12 terms, five white-matter/CSF
principal components, and linear and quadratic trends. The global signal is
included only in the GSR strategies \cite{PCPdocumentationRevision}.
We use filtered ROI time series from \texttt{filt\_global} and
\texttt{filt\_noglobal} for FC, and maps and masks from the corresponding
\texttt{nofilt} strategy for fALFF aggregation. The documented FC filtering band
is 0.01--0.1 Hz; no additional local temporal filter or fixed-length truncation
was applied. Primary time-series lengths range from 116 to 296 volumes.

TR was recovered from inherited site acquisition metadata, including a
FILE\_ID-specific site-variant resolution where necessary. Among the 140 eligible
participants, TR is 1.5 s ($n=19$), 1.667 s ($n=16$), 2.0 s ($n=87$), 2.2 s
($n=1$), or 3.0 s ($n=17$). Consequently, a universal 0.25-Hz Nyquist limit is
inappropriate. Mean FD comes from PCP summary metadata; per-volume censoring
histories are not reconstructed from those summaries.

\textbf{Unresolved provenance:} strategy names do not establish the input branch
that generated the historical fALFF maps. The archived configuration specifies
0.01--0.1 Hz fALFF bounds \cite{PCPconfigRevision}, whereas the submitted text described 0.01--0.08 Hz.
The output mapping identifies a smoothed standard-space fALFF derivative.
The historical implementation, amplitude-versus-power definition, denominator,
filtering order, and smoothing kernel must be verified before this paragraph can be finalized.
We do not infer these facts from a current C-PAC release or from equal maps
across filtered/unfiltered strategy folders.

\subsection*{Coverage-defined nodes and cortical assignments}
CC200 parcel labels were resampled by nearest-neighbor interpolation to each
fALFF map grid. Regional fALFF is the arithmetic mean of finite map values within
the parcel and the downloaded functional mask; coverage is the fraction of parcel
voxels meeting these criteria. A retained node must have a nonconstant downloaded
time series and at least 50\% finite in-mask support in both pipelines.
The primary set contains the 106 nodes satisfying this rule in every eligible
participant. All 140 participants therefore have the same primary nodes.
This is a coverage-limited network rather than an unchanged whole-brain CC200 graph.
The threshold is a methodological choice, not a universal validity standard.

Two exploratory designs assess scope: 181 nodes qualifying in at least 95\% of
eligible participants, followed by requiring every selected node to qualify in
each retained participant (104 participants, 48 ASD/56 controls); and all 200 nodes
with the same participant-level rule (69 participants, 31/38). Medians and all
features are recomputed within each node set. These designs change both nodes and
participants, precluding attribution of between-design differences to anatomy alone.

Cortical network labels use the liberal Yeo seven-network atlas \cite{Yeo2011}.
Parcels require at least 50\% Yeo cortical support and are then assigned by maximum
Dice overlap. Other nodes are unassigned for cortical blocks but retained globally.
The 106-node set includes 91 assigned and 15 unassigned nodes. Only one assigned
node is limbic; within-limbic summaries are structurally unavailable, not null tests.
The atlas mirror's inconsistent float32 bit-depth header was normalized after
payload and label validation; voxel values and spatial transforms were preserved.

\subsection*{Hamiltonian and complementary features}
For participant $k$, let $a_{ik}$ be the regional downloaded fALFF mean and
$m_k$ the median over the selected nodes. We define
\begin{equation}
s_{ik}=\begin{cases}+1,&a_{ik}\geq m_k,\\-1,&a_{ik}<m_k.\end{cases}
\end{equation}
These are relative mathematical labels: the same fALFF value can receive opposite
labels in two participants. An above-median label does not imply a higher mean
neuronal firing rate. The median split also removes absolute between-person
fALFF scaling from the state definition.

Pearson correlations of the downloaded filtered regional series are clipped to
$[-0.9999,0.9999]$, Fisher transformed, and assigned zero diagonal. With these
weights $z_{ij,k}$, the primary statistic is
\begin{equation}
H_k=-\sum_{i<j}z_{ij,k}s_{ik}s_{jk}.\label{eq:H}
\end{equation}
In Eq.~\ref{eq:H}, lower values indicate greater alignment under this definition,
not lower biological
energy or demonstrated dynamical stability. The four pairwise sign/state classes
are illustrated in Fig.~\ref{fig:configuration}. Negative FC is a statistical association,
not a direct inhibitory synapse. We retain GSR and no-GSR as paired analyses because
GSR affects signed associations; GSR is not treated as biologically definitive.

Complementary features include four proportions partitioning nonzero edges by
weight sign and nodal-state agreement, and global/local efficiency of the absolute
Fisher-weight graph \cite{latora2001}. Edge lengths are $1/|z_{ij}|$, with zero
weights treated as absent; local efficiency averages the induced-neighborhood
global efficiencies. For each of seven within-network and 21 between-network
blocks we compute mean weight, mean state-weighted weight, and their negative sum.
Edges involving unassigned nodes remain in $H$ but not cortical block sums.
The previous second-eigenvector ``bipolarity'' descriptor is exported only as a
legacy descriptor and is excluded from inference.

\subsection*{Group inference and harmonization}
Primary inference uses untrimmed raw $H$ in ordinary least squares with diagnosis,
site indicators, age, FIQ, and mean FD. We use HC3 standard errors and
$t$-referenced 95\% confidence intervals. Each design excludes sites with fewer
than three retained participants to avoid singleton leverage; the primary model
includes 136 participants (67/69) at nine sites. Design rank is checked.
Holm correction covers the two pipeline diagnosis tests. Secondary whole-network,
within-network, and between-network families use Benjamini--Hochberg correction
across both pipelines within each design/model, excluding unavailable tests and
primary $H$. There are respectively 12, 36, and 126 valid secondary tests in the
primary design. Sensitivity analyses are exploratory rather than further
confirmatory primary families.

ComBat adjusts the valid feature matrix with site as batch, explicit categorical
diagnosis protection, and continuous age, FIQ, and FD protection
\cite{Johnson2007,Fortin:2018neurocomBat}. The same adjusted regression is then fitted.
Failures cannot silently substitute raw values. Uncertainty does not integrate
estimated harmonization parameters. Pooled 1.5-IQR trimming of raw $H$, before
reapplying the site-size rule, is a separate sensitivity; no diagnosis-specific
trimming defines the primary analysis or prediction sample.

Additional reviewer diagnostics summarize site/group $H$, leverage and Cook's
distance, and site partial $R^2$ conditional on biological covariates. Residual
variance comparisons and post-ComBat site summaries are descriptive because the
same sample estimates harmonization. An exploratory raw adjusted analysis excludes
sites with only one diagnosis, with Holm correction over pipelines per design.
No participant is removed on the basis of influence diagnostics.

We report unadjusted Hedges' $g$ with 2,000 within-group participant bootstrap
replicates (seed 101). These intervals are descriptive and do not resample sites.
Sensitivity power uses a two-sample approximation at 80\% power, not a
site/covariate-adjusted power calculation. Exploratory two-one-sided equivalence
tests use bounds of $\pm0.2$ observed pooled SD for the adjusted coefficient;
these are conventional bounds, not a validated clinically negligible effect.

\subsection*{Efficiency associations and randomization checks}
Within-group partial correlations adjust for mean Fisher FC, FD, age, IQ, and site.
Sites becoming singletons within a group are excluded from that group's residual
analysis. Eight tests (two efficiency summaries, two groups, two pipelines) receive
BH correction. Group specificity is tested directly by diagnosis-by-standardized-
global-efficiency interaction with the same controls, Holm corrected over pipelines.
A uniform edge shift $c$ changes $H$ by
\begin{equation}
\Delta H=\frac{n-(\sum_i s_i)^2}{2}c.
\end{equation}
For exactly balanced 106-node states this is $53c$, demonstrating a mathematical
route for global-weight shifts to affect $H$ independently of a diagnostic mechanism.

For each participant and pipeline, 1,000 state permutations hold FC and total state
counts fixed. Unrestricted permutations are compared with permutations within
network/hemisphere blocks (hemisphere defined by parcel-centroid $x<0$ versus
$x\geq0$). We compute participant $z$-scores and lower-tail Monte Carlo probabilities
with a plus-one correction; zero null variance is unavailable. Seeds are
$10000+\mathrm{SUB\_ID}+1$ for GSR and $10000+\mathrm{SUB\_ID}$ for no-GSR.
Block permutations preserve coarse organization, not full spatial autocorrelation
\cite{Markello2021revision}; they are alignment checks rather than independent
validation or diagnosis tests.

To report a non-autistic baseline separately, we summarize these existing null
outputs for all 71 eligible male controls, without energy trimming or the
regression-specific small-site exclusion. Individual lower-tail probabilities
receive an additional exploratory BH correction over all 284 control/pipeline/null
comparisons. Mean energy differences and mean null $z$-scores have descriptive
95\% percentile intervals from 5,000 site-cluster bootstrap replicates
(seed 20261026), conditional on the saved null summaries. These intervals are
not pooled randomization tests and omit Monte Carlo estimation uncertainty.
The signed weighted FC matrix is fixed in these nulls; no edge/weight rewiring
or independent edge-sign reassignment is performed.

\subsection*{Descriptive H--metric complement and original spectral descriptor}
The original descriptive H-versus-feature question is retained as an additional
revision complement. Pearson and Spearman correlations cover all 87 available
nonlegacy features other than H, within each diagnosis and pipeline. BH is
computed separately over 87 tests per group/pipeline/estimator, matching the
original Pearson family convention. An alternative pooled BH over 348 tests per
estimator/version is supplied separately; pooling families is not uniformly more
conservative in every subgroup. Raw eligible samples, raw adjusted-model subsets,
and saved ComBat features on the same model subsets are reported as distinct
versions. These analyses do not change frozen group-inference families.
The spectral bipolarity value from the original second-eigenvector construction
is additionally compared by raw HC3-adjusted diagnosis models with Holm over
two pipeline tests, separately from frozen families. It is an operational
spectral descriptor, not an asserted optimal signed-balance partition.

\subsection*{Clinical and predictive analyses}
The primary design supplies the same 88 available, nonlegacy features to both
pipelines and all clinical/predictive feature analyses, including both efficiencies.
Five ASD outcomes are considered: ADI-R social, verbal, and onset scores, ADOS
total, and ADOS communication. Missing/sentinel scores are excluded per outcome.
Pearson and Spearman feature correlations receive separate BH corrections across
all feature/score/pipeline tests. Adjusted $H$ correlations control age, IQ, FD,
and site, with BH across available score/pipeline tests.

PLS regression uses nested site-disjoint GroupKFold (up to five outer and three
inner folds), selecting one to three components by inner-fold MSE. Imputation,
variance filtering, and PLS scaling are fitted inside training folds. Complete
OOF predictions are compared with each outer training sample's mean outcome.
Exploratory PLSC uses nuisance-residualized and standardized features/outcomes for
ADOS and five-score complete-case blocks, with 1,000 within-site residual
permutations (seed 101), rank-wise singular-value tests, and BH over all reported
latent-variable tests. Its residual exchangeability assumptions and small complete
cases limit interpretation \cite{Krishnan2011}.

% CLASSIFICATION_METHODS_BEGIN
Classification compares logistic regression, RBF SVM, Gaussian naive Bayes, KNN, XGBoost on all 140 primary
participants without ComBat. Deterministic size-balanced site-disjoint folds use
five outer and three inner splits, each required to contain both classes.
Median imputation, variance filtering, standardization, ANOVA feature selection
(25/50/100 percent), and hyperparameter tuning are nested. Logistic and SVM
class weights are balanced. Logistic $C$ values are 0.01/0.1/1/10; SVM $C$ values
are 0.1/1/10 with scale-based gamma. Gaussian naive Bayes uses its default
variance smoothing and training-estimated priors, with feature-selection percentile
tuned inside folds. KNN evaluates 3/5/7/9/11 neighbors and uniform/distance
weights. Positive-class probabilities supply ROC scores for these added models;
logistic/SVM decision-function scores and predictions are reused unchanged.
XGBoost uses CPU histogram trees, seed 42, one thread, log-loss evaluation,
max depth 2/3/4, 100/200/300 estimators, and learning rates .05/.1. The same
nested feature-selection percentiles apply; default scale-pos-weight is retained.
Inner scoring uses balanced accuracy. No model is selected by outer-test accuracy.
Pooled OOF balanced accuracy is accompanied by a 1,000-replicate site bootstrap
conditional on fixed predictions, using the same seed 303 and site order as the
completed Stage 2 models. This omits some training uncertainty. The node mask and
available feature schema use cohort-wide QC, so prediction is exploratory
conditional on this mask, not prospective validation of the entire workflow.
% CLASSIFICATION_METHODS_END

% ML_SELECTION_METHODS_BEGIN
\paragraph*{Exploratory comparison of feature-selection methods.}
After inspecting the baseline outer results, we specified a separate follow-up
using the same 88-feature inputs, participants, and five/three outer/inner site
folds. Imputation, variance filtering, redundancy pruning, scaling, selection,
and classifier fitting are estimated only from the corresponding training data
\cite{scikit-learn}. Redundancy pruning retains the first feature in fixed schema
order when absolute training correlation is at least $1-10^{-10}$; proportional
EnergyBlock/AvgFalffW terms and complementary proportions motivate this step.
The graph, participants and primary group tests are unchanged.

Inner CV jointly chooses all remaining features, or ANOVA, mutual information,
or L1-logistic coefficient-ranked top 5/10/20 features, with K capped at the
available count. Mutual information uses three neighbors and seed 101. The L1
ranker uses balanced liblinear logistic regression, C=1 and seed 101; coefficient
ties are broken by training ANOVA rank and fixed feature order. Zero-coefficient
selections are explicitly recorded; none occurred in the final outer fits.
Non-XGBoost classifier grids match the baseline. The compact follow-up XGBoost
grid uses depth 2/3, estimators 100/200 and learning rates .05/.1, with the same
fixed CPU histogram settings. Thus changes cannot be attributed to selector type
alone. All candidates are selected by inner balanced accuracy, and all ten
model/pipeline results are reported without choosing an outer-test winner.

Paired 1,000-replicate site bootstrap (seed 303) estimates conditional intervals
for the balanced-accuracy change on aligned OOF predictions. These intervals
omit full training uncertainty and do not adjust for repeated exploration.
Because baseline outer results had already been inspected, this follow-up is
exploratory rather than independent confirmation of improved generalization.
External evaluation would be needed to substantiate a predictive gain.
% ML_SELECTION_METHODS_END

\section*{Data and code availability}
ABIDE I and PCP derivatives are publicly available from their source repositories.
The analysis includes frozen participant/ROI manifests, input hashes, derivative
provenance records, and executable revision scripts.
\textbf{Author action required:} insert the verified public archive/repository
identifier for the exact revision release and its permitted underlying data.
Existing Kaggle datasets supply inputs but do not by themselves identify the
complete final code/results release.

\section*{Acknowledgments}
We thank ABIDE, its contributing sites and participants, and the Preprocessed
Connectomes Project for making these data and derivatives available.

\nolinenumbers
\bibliography{Bibliography_revision}

\section*{Figure legends}
\begin{figure}[htbp]
  \centering
  \includegraphics[width=.85\textwidth]{figures/participant_flow.pdf}
  \caption{Participant flow for the revised primary analysis. Demographic eligibility
  precedes selection of a common coverage-defined node set. Four participants are
  excluded only from adjusted group models by the minimum-site-size rule;
  classification retains all 140 eligible participants.}\label{fig:flow}
\end{figure}
\begin{figure}[htbp]
  \centering
  \includegraphics[width=.96\textwidth]{figures/control_null_empirical.pdf}
  \caption{Male-control empirical energy against each participant's own null mean.
  Columns show GSR/no-GSR; rows show unrestricted and network/hemisphere-blocked
  state permutations. Every panel contains the same 71 eligible male controls.
  Lines connect each control's empirical $H$ with the mean of 1,000 null
  configurations; black diamonds indicate across-control means. The complete
  signed weighted FC matrix is fixed and states are reassigned. These are not
  edge-rewired networks or full spatial-autocorrelation nulls.}
  \label{fig:controlnull}
\end{figure}
\begin{figure}[htbp]
  \centering
  \includegraphics[width=\textwidth]{figures/primary_results.pdf}
  \caption{Primary group coefficients and exploratory classification. Left: raw
  adjusted diagnosis coefficients with HC3 95\% intervals ($n=136$). Right:
  pooled site-held-out balanced accuracy ($n=140$), with conditional site-bootstrap
  95\% intervals. The vertical references are zero and chance accuracy, respectively.
  Prediction intervals omit some model-training uncertainty.}\label{fig:primary}
\end{figure}
\begin{figure}[htbp]
  \centering
  \includegraphics[width=\textwidth]{figures/global_feature_distributions.pdf}
  \caption{All seven nonlegacy global features in the untrimmed 106-node feature
  cohort: $H$, two absolute-weight graph efficiencies, and four pairwise
  configuration proportions. Points show 69 ASD and 71 control participants per
  pipeline; boxes show quartiles and medians. These are raw descriptive
  distributions; adjusted group inference uses 136 participants and is reported
  in the corresponding model table. No feature is selected by its $p$-value.}
  \label{fig:globalfeatures}
\end{figure}
\begin{figure}[htbp]
  \centering
  \includegraphics[width=\textwidth]{figures/intra_network_effects.pdf}
  \caption{Within-network ASD--control comparisons for all seven named cortical
  systems and all three original block measures in both pipelines. Colors show
  diagnosis coefficient divided by HC3 SE. Gray limbic cells are unavailable,
  not nonsignificant tests. No valid test survives its specified BH correction.}
  \label{fig:intraeffects}
\end{figure}
\begin{figure}[htbp]
  \centering
  \includegraphics[width=.93\textwidth,height=.68\textheight,keepaspectratio]{figures/inter_network_effects.pdf}
  \caption{Between-network ASD--control comparisons for all 21 network pairs
  and all three original block measures. Colors show HC3 t statistics, not
  corrected significance. Complete estimates and correction scopes accompany
  the figure in S2. No comparison survives the frozen family correction.}
  \label{fig:intereffects}
\end{figure}
\begin{figure}[htbp]
  \centering
  \includegraphics[width=\textwidth,height=.68\textheight,keepaspectratio]{figures/H_correlates_raw.pdf}
  \caption{Top 20 descriptive correlates of H within each diagnosis/pipeline,
  restoring the original within-group analysis. Ranking is by absolute Pearson
  correlation among all 87 available other features. Asterisks indicate BH
  q$<.05$ over those 87 tests per panel. Red/blue bars indicate positive/negative
  correlations. Abbreviations: VIS, visual; SM, somatomotor; DA, dorsal attention;
  SV, salience/ventral attention; LIM, limbic; FP, frontoparietal; DMN, default;
  W, mean FC; F, state-weighted FC; E, block energy. Shared terms and proportional
  F/E pairs make many associations mathematically dependent. All tests and
  matched raw/ComBat versions are in S2; these are not ASD--control contrasts.}
  \label{fig:Hcorrelates}
\end{figure}
\begin{figure}[htbp]
  \centering
  \includegraphics[width=\textwidth]{figures/pair_configurations.png}
  \caption{Pairwise configurations of relative nodal states and FC signs. Magenta
  and teal nodes denote $+1$ and $-1$ mathematical states, respectively; green and
  magenta links denote positive and negative FC. Same-state positive edges
  (agreement) and opposite-state negative edges (disagreement) both lower $H$;
  the other two configurations raise it. Nodes do not denote neuronal firing
  states. This conceptual illustration is retained from the submitted article.}
  \label{fig:configuration}
\end{figure}

\section*{Supporting information legends}
\paragraph{S1 Supporting Information.} Additional methods, cohort/site summaries,
model diagnostics, node-set sensitivities, H--metric complements, control-only null
comparisons, adjusted efficiency plots, clinical validation, and complete OOF
ROC/confusion plots, provided as a separate PDF.
\paragraph{S2 Supporting Information.} Machine-readable supporting tables and
manifests with a file index and provenance hashes, provided as a separate ZIP.
This local draft archive requires author review of sharing terms before deposit.

\end{document}
